\problem{MEC-0037}{37}
(IJSO รอบที่ 1 2552) อนุภาคหนึ่งถูกโยนขึ้นในแนวดิ่งให้เคลื่อนที่ภายใต้สนามโน้มถ่วง $g$ กราฟความเร็ว--เวลาของอนุภาค เป็นดังรูปข้างล่าง (ทิศขึ้นเป็นทิศที่แทนด้วยเครื่องหมายบวก) ตอนขาลง เมื่ออนุภาคอยู่ที่ตำแหน่งครึ่งหนึ่งของความสูงสูงสุด เวลาผ่านไปเท่าใดนับจากเริ่มโยน
\begin{center}\begin{tikzpicture}[>=Stealth]\draw[->] (0,-1.1)--(0,2.4) node[above]{ความเร็ว};\draw[->] (-.3,0)--(5,0) node[right]{เวลา};\draw[thick] (0,1.6)--(4.4,-1.1);\node[below] at (2.61,0){$T$};\end{tikzpicture}\end{center}
\par\smallskip\par\smallskip\begin{tabular}{@{}p{.48\linewidth}p{.48\linewidth}@{}}1. $\dfrac{T}{2}$ & 2. $T\left(1+\dfrac{1}{\sqrt{2}}\right)$ \\ 3. $\dfrac{3T}{2}$ & 4. $T\left(2-\dfrac{1}{\sqrt{2}}\right)$\end{tabular}
